% ===================================================================
%  CUADRO N.º 3 — Infancia y Juventud.
%  Traduction 2026 d'apres texto/io, controlee sur texto/fr et
%  texto/en. Consigne : voir texto/es/10-cuadro-01.tex.
% ===================================================================

%%K t03-apar-1 apar
\VUsaut{3.0mm}
\VUtitre{15.0pt}{60}{CUADRO N.º 3}
\VUsaut{1.6mm}
\VUtitre{11.4pt}{0}{Infancia y Juventud.}
\VUsaut{3.4mm}

%%K t03-apar-2 apar
(Los cuadros 3 y 4 están dispuestos de modo que presentan, en cuatro
\VUgras{escenas familiares}, el vocabulario esencial del \VUgras{tiempo},
la \VUgras{edad}, la \VUgras{familia}, la \VUgras{ropa} y la \VUgras{salud}.)
\VUsaut{4.0mm}

%%K t03-tit-1 sub
\VUcentre{12.0pt}{0}{\textit{Primera escena.}}
\VUsaut{3.0mm}

%%K t03-tit-2 sub
\VUpk{10.2pt}{40}{Un bautizo en el pueblo.}
\VUsaut{2.4mm}

%%K t03-tit-3 sub
\VUpk{10.2pt}{40}{La primavera.}
\VUsaut{3.0mm}

%%K t03-c1-01-1 p
1. --- Estamos en \VUgras{primavera}, en el mes de \VUgras{marzo},
\VUgras{abril} o \VUgras{mayo}, en un día de la \VUgras{semana} ---
\VUgras{lunes}, \VUgras{martes} o \VUgras{miércoles}. Son las
\VUgras{diez} de la mañana.

%%K t03-c1-02-1 p
2. --- Hace buen \VUgras{tiempo}, el \VUgras{aire} es puro. El sol brilla.
Unas \VUgras{nubecillas} \textsuperscript{(2)} suben por el \VUgras{cielo}
\textsuperscript{(1)} azul.

%%K t03-c1-03-1 p
3. --- Los \VUgras{árboles} apenas tienen \VUgras{hojas}. Los delgados
\VUgras{álamos} \textsuperscript{(3)} y el alto \VUgras{plátano} \textsuperscript{(17)} no
llevan todavía más que brotes.

%%K t03-c1-04-1 p
4. --- Las \VUgras{golondrinas} \textsuperscript{(4)} han vuelto y revolotean con
alegres chillidos alrededor de la \VUgras{flecha} \textsuperscript{(7)} del
\VUgras{campanario} \textsuperscript{(5)} que corona la \VUgras{iglesia}
\textsuperscript{(6)} del pueblo.

%%K t03-tit-4 sub
\VUsaut{3.4mm}
\VUpk{10.2pt}{40}{La Iglesia.}
\VUsaut{3.0mm}

%%K t03-c1-05-1 p
5. --- Muy sencilla es esta iglesia con su gran \VUgras{portada}
\textsuperscript{(13)} y su grueso \VUgras{reloj} \textsuperscript{(8)}. La
\VUgras{aguja pequeña} \textsuperscript{(9)} está sobre la cifra X: marca las
\VUgras{horas}. La \VUgras{grande} \textsuperscript{(10)} está sobre las XII
(mediodía); marca los \VUgras{minutos}.

%%K t03-c1-06-1 p
6. --- En el campanario, las \VUgras{campanas} \textsuperscript{(11)} suenan
alegremente. En lo alto del tejado se alza una pequeña \VUgras{cruz}
\textsuperscript{(12)} de piedra.

%%K t03-c1-07-1 p
7. --- La iglesia no tiene solo una gran \VUgras{portada} \textsuperscript{(13)},
tiene también una puertecilla lateral. Por esta puerta el señor
\VUgras{cura} \textsuperscript{(15)}, vestido con su \VUgras{sotana}, sale de la
\VUgras{sacristía} \textsuperscript{(14)} y se dirige a la \VUgras{casa parroquial}
\textsuperscript{(19)}.

%%K t03-c1-08-1 p
8. --- Cerca de él, he aquí la \VUgras{lechera} \textsuperscript{(24)} con su
\VUgras{burro} \textsuperscript{(23)}. Vuelve de la ciudad, adonde ha llevado su
buena leche para los niños.

%%K t03-tit-5 sub
\VUsaut{4.0mm}
\VUpk{10.2pt}{40}{El huerto de frutales.}
\VUsaut{3.4mm}

%%K t03-c1-09-1 p
9. --- El señor Aliborón (el burro) está muy contento de esta carrera
matinal. Aspira con deleite el aire embalsamado por el perfume de las
\VUgras{flores} \textsuperscript{(20)} que cubren los \VUgras{árboles frutales}
\textsuperscript{(18)} del \VUgras{huerto} vecino. Todos rosados y blancos están
esos \VUgras{melocotoneros} \textsuperscript{(18)} y esos \VUgras{albaricoqueros}
\textsuperscript{(19)}, que prometen una buena cosecha.

%%K t03-c1-10-1 p
10. --- Mientras tanto, las pequeñas \VUgras{abejas} \textsuperscript{(22)} de la
\VUgras{colmena} \textsuperscript{(21)} van allí a libar su dulce \VUgras{miel},
zumbando.

%%K t03-c1-11-1 p
11. --- Pero ¿qué pasa? He aquí a los niños del pueblo, que llegan
corriendo y hacen huir a los \VUgras{gansos} \textsuperscript{(25)} asustados. Es el
cortejo que sale de la iglesia después del \VUgras{bautizo} del hijo del
notario.

%%K t03-c1-12-1 p
12. --- El pequeño Jacobo, en su \VUgras{cuna} \textsuperscript{(26)}, se despierta
con el alegre son de las campanas, y su hermana Magdalena, que también
quiere ver el cortejo, le pide a su madre que venga a peinarla y a
vestirla en el umbral de la casa.

%%K t03-c1-13-1 p
13. --- La madre accede. Se pone su \VUgras{pañuelo de cabeza}
\textsuperscript{(28)}, su \VUgras{corpiño} \textsuperscript{(29)} y su \VUgras{delantal}
\textsuperscript{(30)}, y viene a sentarse en una sillita delante de la puerta.
Magdalena solo lleva sus \VUgras{enaguas} \textsuperscript{(31)} y su \VUgras{corsé}
\textsuperscript{(32)}.

%%K t03-c1-14-1 p
14. --- ¡Qué feliz es! Se olvida de sus flores preferidas, sus
\VUgras{claveles} \textsuperscript{(34)}, donde se posan las \VUgras{mariposas}
\textsuperscript{(35)}, sus \VUgras{tulipanes} \textsuperscript{(36)}, sus
\VUgras{muguetes} \textsuperscript{(37)} en \VUgras{macetas} \textsuperscript{(38)} sobre el
\VUgras{muro} \textsuperscript{(33)}, y sus \VUgras{rosas} \textsuperscript{(39)}, que forman
un seto tan hermoso. Contempla los ricos vestidos de las señoras del
cortejo.

%%K t03-c1-15-1 p
15. --- Los chicos del pueblo no hacen mucho caso de los vestidos. Se
precipitan y se empujan para conseguir \VUgras{caramelos} \textsuperscript{(55)} o
\VUgras{monedas} \textsuperscript{(52)}. Uno de ellos llora \textsuperscript{(40)}.

%%K t03-c1-16-1 p
16. --- Otro ha pisado la pata del \VUgras{perro} \textsuperscript{(42)}, que huye
chillando y hace huir a la \VUgras{gallina} \textsuperscript{(43)} y a sus
\VUgras{pollitos} \textsuperscript{(44)}.

%%K t03-tit-6 sub
\VUsaut{5.0mm}
\VUpk{10.2pt}{40}{El cortejo del bautizo.}
\VUsaut{4.0mm}

%%K t03-c1-17-1 p
17. --- Delante del cortejo camina el \VUgras{ama} \textsuperscript{(45)}. Lleva una
hermosa \VUgras{cofia} \textsuperscript{(46)} con cintas largas. Lleva en brazos al
\VUgras{recién nacido} \textsuperscript{(47)}, todo vestido de \VUgras{encajes}
\textsuperscript{(48)}.

%%K t03-c1-18-1 p
18. --- Detrás del ama vienen el \VUgras{padrino} \textsuperscript{(49)} y la
\VUgras{madrina} \textsuperscript{(53)}. Reparten los caramelos y las monedas. El
padrino lleva \VUgras{guantes} \textsuperscript{(51)} y \VUgras{sombrero de copa}
\textsuperscript{(50)}. La madrina lleva en la mano un hermoso \VUgras{ramo}
\textsuperscript{(54)}.

%%K t03-c1-19-1 p
19. --- Detrás de ellos vemos al \VUgras{tío} \textsuperscript{(56)} y a la
\VUgras{tía} \textsuperscript{(58)} del niño. La tía lleva un elegante
\VUgras{sombrero} \textsuperscript{(59)}, el tío una \VUgras{levita} \textsuperscript{(57)}
negra y un chaleco blanco. Vienen luego el \VUgras{primo} \textsuperscript{(60)} y la
\VUgras{prima} \textsuperscript{(61)}. Esta lleva de la mano a la \VUgras{hermana}
\textsuperscript{(62)} del recién nacido, la pequeña Berta.

%%K t03-c1-20-1 p
20. --- El otro \VUgras{hermano} \textsuperscript{(63)} de Berta camina detrás. Da la
mano a una \VUgras{parienta} \textsuperscript{(64)}. Los \VUgras{amigos}
\textsuperscript{(65)} de la familia vienen después.

%%K t03-tit-7 sub
\VUsaut{4.6mm}
\VUpk{10.2pt}{40}{Los que miran.}
\VUsaut{3.4mm}

%%K t03-c1-21-1 p
21. --- Cerca del cortejo, he aquí a un \VUgras{mendigo} \textsuperscript{(66)}
apoyado en su \VUgras{muleta} \textsuperscript{(67)}. Pide limosna.

%%K t03-c1-22-1 p
22. --- Al otro lado hay un niño muy \VUgras{cortés} \textsuperscript{(68)}. Saluda y
da los \VUgras{«buenos días»} a las personas que conoce.

%%K t03-c1-23-1 p
23. --- Los del pueblo han salido de sus casas para ver pasar el cortejo
del bautizo. Vemos a un \VUgras{campesino} \textsuperscript{(69)} con su
\VUgras{gorro de algodón} y a su lado a una \VUgras{campesina}
\textsuperscript{(70)}. Luego a una \VUgras{anciana} \textsuperscript{(71)} apoyada en su
\VUgras{bastón} \textsuperscript{(72)}. Por último, al \VUgras{molinero}
\textsuperscript{(73)} con su ancho \VUgras{sombrero de fieltro} \textsuperscript{(74)}, y a
una campesina joven calzada con \VUgras{zuecos} \textsuperscript{(75)}, que enseña a
andar a su niño.

%%K t03-c1-24-1 p
24. --- Es una escena muy divertida.
\VUsaut{4.0mm}

%%K t03-tit-8 sub
\VUcentre{12.0pt}{0}{\textit{Segunda escena.}}
\VUsaut{3.0mm}

%%K t03-tit-9 sub
\VUpk{10.2pt}{40}{Una fiesta pública.}
\VUsaut{2.4mm}

%%K t03-tit-10 sub
\VUpk{10.2pt}{40}{Juegos náuticos y regatas.}
\VUsaut{3.0mm}

%%K t03-c2-01-1 p
1. --- Ha llegado el \VUgras{verano}. El sol ardiente del mes de
\VUgras{junio} (julio o \VUgras{agosto}) invita a los jóvenes a bañarse y a
salir en barca.

%%K t03-c2-02-1 p
2. --- En la orilla derecha del río, los remeros se desnudan para tomar
parte en los juegos náuticos y en las regatas.

%%K t03-c2-03-1 p
3. --- Uno de ellos se quita los \VUgras{zapatos} \textsuperscript{(1)}; los
desata (los desabrocha). Sus dos compañeros ya están listos. No llevan
más que su \VUgras{jersey} \textsuperscript{(2)} y su \VUgras{bañador}
\textsuperscript{(3)}. Charlan mientras esperan a sus amigos. Hablan de la feria y
de la fiesta que se celebra en la otra orilla.

%%K t03-c2-04-1 p
4. --- En el suelo, junto a ellos, yace la \VUgras{ropa} de sus compañeros:
un \VUgras{par} de \VUgras{botines} \textsuperscript{(4)}, \VUgras{zapatos}
\textsuperscript{(5)}, \VUgras{calcetines} \textsuperscript{(6)}, sombreros de
\VUgras{fieltro} \textsuperscript{(7)} y de \VUgras{paja} \textsuperscript{(8)} y una
\VUgras{corbata} \textsuperscript{(9)}.

%%K t03-c2-05-1 p
5. --- Uno de los jóvenes ha doblado con cuidado su \VUgras{chaqueta}
\textsuperscript{(10)}. La pone sobre una toalla, en la que su vecino envolverá
enseguida los dos \VUgras{trajes}.

%%K t03-c2-06-1 p
6. --- Un sexto muchacho va con retraso. Todavía lleva su \VUgras{gorra}
\textsuperscript{(11)} en la cabeza. No se ha quitado ni la \VUgras{camisa}
\textsuperscript{(12)}, ni el \VUgras{cuello} \textsuperscript{(13)}, ni los
\VUgras{puños} \textsuperscript{(14)}. Tiene en la mano su \VUgras{chaleco}
\textsuperscript{(17)} y lleva todavía sus \VUgras{pantalones} \textsuperscript{(16)} y sus
\VUgras{tirantes} \textsuperscript{(15)}. Seguro que no estará listo para la
próxima carrera.

%%K t03-c2-07-1 p
7. --- Cerca de los jóvenes, un sendero bordeado de \VUgras{cardos}
\textsuperscript{(18)} baja a la orilla del río, donde el tío Jacobo prepara,
entre las \VUgras{cañas} \textsuperscript{(19)}, los \VUgras{fuegos artificiales}
\textsuperscript{(20)} que se dispararán por la noche.

%%K t03-tit-11 sub
\VUsaut{4.4mm}
\VUpk{10.2pt}{40}{La carrera.}
\VUsaut{3.4mm}

%%K t03-c2-08-1 p
8. --- En el río, algunas señoras que se resguardan bajo sus sombrillas
pasean en \VUgras{barca} \textsuperscript{(21)}. Siguen la carrera, que es muy
interesante. En cada \VUgras{yola} \textsuperscript{(22)}, cuatro \VUgras{remeros}
\textsuperscript{(24)} reman con todas sus fuerzas, mientras el \VUgras{timonel}
\textsuperscript{(26)} sostiene el \VUgras{timón} \textsuperscript{(25)} y dirige la frágil
embarcación. Los \VUgras{remos} \textsuperscript{(23)} golpean el agua a compás, y
los ligeros botes navegan a una velocidad vertiginosa.

%%K t03-c2-09-1 p
9. --- Unos cuantos jóvenes se disputan, junto al pilar del puente, el
premio de la \VUgras{cucaña horizontal} \textsuperscript{(28)}. Uno de ellos avanza
con cautela por el palo flexible y resbaladizo para alcanzar la pequeña
\VUgras{bandera} \textsuperscript{(29)} sujeta en el extremo. Su desdichado
\VUgras{competidor} \textsuperscript{(30)} ha caído al agua. Nada para llegar a la
orilla.

%%K t03-c2-10-1 p
10. --- Unos muchachos mayores \VUgras{justan en barca} \textsuperscript{(32)} junto
al \VUgras{muelle} \textsuperscript{(33)}. A todos estos juegos asiste un
\VUgras{público} \textsuperscript{(31)} numeroso, apoyado en el \VUgras{pretil} del
\VUgras{puente} \textsuperscript{(27)} y apiñado en la \VUgras{orilla}. Algunos
espectadores, sin embargo, se vuelven para seguir el juego de la
\VUgras{cucaña} \textsuperscript{(45)} o para admirar el \VUgras{cortejo municipal}
que acude al \VUgras{reparto} de los premios.

%%K t03-c2-11-1 p
11. --- Primero, he aquí a unos \VUgras{golfillos} \textsuperscript{(35)} que van
delante de los \VUgras{músicos} \textsuperscript{(36)}; luego vienen el
\VUgras{alcalde} \textsuperscript{(37)}, los tenientes de alcalde y el
\VUgras{ayuntamiento} de la pequeña ciudad. Cierran la marcha los
\VUgras{bomberos} \textsuperscript{(38)}.

%%K t03-tit-12 sub
\VUsaut{5.0mm}
\VUpk{10.2pt}{40}{La feria.}
\VUsaut{3.6mm}

%%K t03-c2-12-1 p
12. --- Hay mucha gente en el \VUgras{real de la feria} \textsuperscript{(39)}. Se
viene a ver las distintas \VUgras{barracas}: los \VUgras{caballitos de
madera} \textsuperscript{(40)}, el \VUgras{circo} \textsuperscript{(41)}, donde la función
ya ha empezado; el \VUgras{baile público} \textsuperscript{(42)}, donde los jóvenes y
las jóvenes de la ciudad gozan del placer de \VUgras{bailar}.

%%K t03-c2-13-1 p
13. --- Al fondo vemos la \VUgras{plaza de toros} \textsuperscript{(44)}, donde el
valiente \VUgras{matador} español lucha contra el \VUgras{toro}
\textsuperscript{(45)} furioso; luego la \VUgras{montaña rusa} \textsuperscript{(49)} y la
\VUgras{barraca de tiro} \textsuperscript{(46)}, donde se ejercitan en manejar las
\VUgras{armas de fuego} y en tirar al \VUgras{blanco}.

%%K t03-c2-14-1 p
14. --- Cerca de allí, he aquí el \VUgras{teatro de feria} \textsuperscript{(47)}, en
el que se representan la \VUgras{comedia} y el \VUgras{drama}. Muy cerca está
la barraca del \VUgras{gigante} \textsuperscript{(48)} y del \VUgras{enano}. La
\VUgras{estatura} del primero es de dos metros quince; el segundo apenas es
tan alto como un niño.

%%K t03-c2-15-1 p
15. --- Por último, he aquí la gran \VUgras{colección de fieras}
\textsuperscript{(50)} con sus \VUgras{animales salvajes}, su \VUgras{tigre}
\textsuperscript{(51)} de poderosas \VUgras{garras}, su \VUgras{león} \textsuperscript{(52)}
de espesa \VUgras{melena}, sus dos \VUgras{jirafas} \textsuperscript{(53)} y su
\VUgras{elefante} \textsuperscript{(54)}, que usa con tanta maña su \VUgras{trompa}.

%%K t03-c2-16-1 p
16. --- El \VUgras{domador} \textsuperscript{(55)} que amaestra a estos animales está
a la puerta de la barraca.
\VUsaut{6.0mm}
\VUfilet{22mm}
